\documentclass[a4paper,11pt]{article}
\usepackage[margin=25mm]{geometry}
\usepackage{amsmath,amssymb,amsthm,booktabs,xcolor}
\usepackage[unicode,colorlinks=true,linkcolor=blue!45!black,urlcolor=blue!45!black,citecolor=blue!45!black]{hyperref}
\usepackage{enumitem}
\usepackage{xurl}
\setlist{nosep}
\setlength{\emergencystretch}{2em}
\newtheorem{theorem}{Theorem}[section]
\newtheorem{proposition}[theorem]{Proposition}
\newtheorem{lemma}[theorem]{Lemma}
\newtheorem{corollary}[theorem]{Corollary}
\newtheorem{remark}[theorem]{Remark}
\numberwithin{equation}{section}
\newcommand{\E}{\mathbb E}
\newcommand{\Var}{\operatorname{Var}}
\newcommand{\Cov}{\operatorname{Cov}}
\newcommand{\Pois}{\operatorname{Pois}}
\newcommand{\given}{\,\middle|\,}
\title{Supplementary appendix: key lemmas for the macroscopic conditional law\\\large Self-contained proofs corresponding to Propositions 8 and 10 of the main text}
\author{Min Wu}
\begin{document}
\maketitle

\setcounter{section}{5}
\renewcommand{\thesection}{\Alph{section}}
\section{Key lemmas for the macroscopic conditional law}\label{app:macro}

\subsection{Scope, notation, and existing inputs}

Throughout we fix $k\ge3$, work in the Plain model, and place the load in the moderate deviation regime:
\begin{equation}
\frac dM=\alpha_c+\frac{r_M}{\sqrt M},\qquad r_M\to-\infty,\qquad |r_M|\le C\sqrt{\log M}.\label{F:range}
\end{equation}
Write $t_0=-g_kr_M\to\infty$, where $g_k=ka_ce^{-x_c}/\sigma_k$; $p_M=\Phi(-t_0)$;
$\lambda_M=\binom d2/\binom Mk$. The macroscopic event is
$A=\{\text{the number of edges of the actual 2-core}>\varepsilon_0M\}$. The sign factor is
$\gamma_M=(d\rho_M^2-1)/(d-1)$; the all-positive-sign measurement is denoted by $X_+$, and the
measurement under an arbitrary fixed sign composition is denoted by $X_s$. $\ell(t)=\phi(t)/\Phi(-t)$
denotes the left-tail Mills ratio, namely $A(-t)$ of the main text.

This appendix proves four groups of conclusions, related to the main text as follows.
\begin{enumerate}[label=(\arabic*)]
\item The macroscopic conditional CDF of the all-positive-sign measurement:
$\sup_x|\Pr(X_+\le x\mid A)-\Phi_2(x,-t_0;\eta_k)/\Phi(-t_0)|\to0$ (Proposition~\ref{F:prop:cdf}).
\item The conditional moments of the all-positive-sign measurement and the conditional mean of an arbitrary
sign composition (Proposition~\ref{F:prop:moments}).
\item The macroscopic conditional CDF and conditional variance of an arbitrary sign composition, namely the
first two formulas of Proposition 8 of the main text (Proposition~\ref{F:prop:sign}).
\item The control of the three weighted error terms, namely the weighted-mean formula of Proposition 10 of
the main text (Proposition~\ref{F:prop:wm}).
\end{enumerate}

We directly use the scalar inputs already established in Appendix C (the proof of Proposition 7 of the main
text). For ease of reference, they are collected here.

\begin{enumerate}[label=(C\arabic*)]
\item There is a good event $G_M$ whose complement has super-polynomially small probability. On $G_M$, up to
the stopping time, all histogram errors are at most $C\vartheta_M$, where
$\vartheta_M=\sqrt M(\log M)^6$.
\item In the configuration model, the initial Doob increments and the dynamic increments assemble into a score
martingale $S$; its increments are at most $C/\sqrt M$, the bracket has a deterministic upper bound, and,
except for a super-polynomially small exceptional probability, it lies within $1\pm v_M$, where
$v_M=C\{M^{-1/8}+(\log M)^8/\sqrt M\}$.
\item There is a bandwidth sandwich between the macroscopic event and the score event:
$A\,\triangle\,\{S\le-b_M\}\subseteq\mathcal B_M\cup\{|S+b_M|\le C\delta_M\}$,
where $\mathcal B_M$ is super-polynomially small, the deterministic bandwidth $\delta_M\to0$ satisfies
$b_M\delta_M\to0$, and $b_M=t_0+O((1+r_M^2)/\sqrt M)$.
\item The relative asymptotics of the macroscopic probability: $\Pr(A)=p_M(1+o(1))$; the same holds with $l$
fixed special edges, and it can be transferred to the Plain model via relative Bonferroni.
\end{enumerate}

``Super-polynomially small'' means $o(M^{-b})$ for every fixed $b>0$. By \eqref{F:range}, $p_M$ is no smaller
than a fixed negative power divided by a polylogarithm, so a super-polynomially small quantity divided by $p_M$
remains super-polynomially small. This is used repeatedly below.

\subsection{Choice of the macroscopic cutoff}

Proposition 7 of the main text holds for any $0<\varepsilon_0<\alpha_cq_c^k$. The two-regime reduction also
requires the small-core counts to cover all core sizes below the macroscopic layer, so this appendix fixes a
sufficiently small cutoff.

\begin{lemma}[Choice of the cutoff]\label{F:lem:cutoff}
Take $\varepsilon_*>0$ such that the small-core summation of Eq.~(64) of Appendix B holds for $r\in\{0,2\}$ on
$3\le e\le\varepsilon_*M$, and such that $\varepsilon_*<\alpha_cq_c^k/2$. Fix $0<\varepsilon_0\le\varepsilon_*$.
Let $D$ be the event that a repeated edge pair exists, and let $R=F\setminus(A\cup D)$. Then under
\eqref{F:range}
\begin{equation}
\Pr(R)=O(M^{-b_k})=o(\lambda_M),\qquad \Pr(A\cap D)=o(\lambda_M).\label{F:cutoff}
\end{equation}
\end{lemma}

\begin{proof}
On $R$, the actual core has between $3$ and $\varepsilon_0M$ edges. The first-moment summation of Appendix B
requires only $d\le AM$, independently of the distance from the threshold, so $\Pr(R)=O(M^{-b_k})$; moreover
$b_k-(k-2)\ge1$, so this is $o(\lambda_M)$.

For the intersection, first consider a fixed window. Suppose $d/M=\alpha_c+r/\sqrt M+o(M^{-1/2})$ with $r$
fixed, plant a pair of identical uniform $k$-subsets, and let the remaining edges be independent. The finite
planting lemma of Section 5.5 allows the special edges to be mutually correlated, so the precritical score
still has the same Gaussian limit; the two repeated edges are themselves legitimate $k$-subsets, the
probability of having no bad edge among the remaining regular edges has a positive limit, and the
inclusion--exclusion transfer proceeds as usual. The planted repeated edge makes final failure inevitable, so
here only the macroscopic event is transferred: when the score lies below a fixed negative gap, the process
stops with $(\alpha_cq_c^k+o(1))M$ edges remaining, which belongs to $A$; when the score lies above a fixed
positive gap, the concentration before stopping and the local freezing rule out early stopping until fewer
than $\varepsilon_0M/2$ edges remain, which does not belong to $A$. Removing the safety gap gives
$Q_r(A)\to\Phi(g_kr)$, where $Q_r$ is the distribution after planting.

Return to \eqref{F:range}. For any fixed $r_0<0$, eventually $d\le\lfloor\alpha_cM+r_0\sqrt M\rfloor$.
Couple by adding edges using the same sequence of regular edges; the macroscopic event is monotone in the
addition of edges, so $\limsup_MQ_M(A)\le\Phi(g_kr_0)$. Letting $r_0\to-\infty$ gives $Q_M(A)\to0$. The
density of the random-label planting distribution relative to the original distribution is $N/\lambda_M$,
where $N$ is the number of repeated edge pairs, so
\[
\Pr(A\cap D)\le\E[N\mathbf1_A]=\lambda_MQ_M(A)=o(\lambda_M).\qedhere
\]
\end{proof}

The probability asymptotics of Proposition 7 of the main text hold for any $\varepsilon_0<\alpha_cq_c^k$, so
the macroscopic probability itself is independent of the choice of the cutoff. Proposition 8, Proposition 10,
and the conclusions of this appendix conditioned on $A$ all refer to the cutoff chosen in
Lemma~\ref{F:lem:cutoff}. When incorporated into the main text, the hypothesis of Proposition 8 is
correspondingly changed to ``take $\varepsilon_0$ sufficiently small''.

\subsection{All-positive signs: the two-dimensional initial Doob decomposition}

We first work in the configuration model. Let $n_e=kd$ be the total number of endpoints, $\mu=n_e/M$, the
endpoints $V_1,\ldots,V_{n_e}$ independent uniform, and $D_v$ the occupancy count. Define
\begin{equation}
K=\sum_v(D_v)_2,\qquad Y_K=\frac{K-\E K}{\sqrt2\mu\sqrt M},\qquad H=\sum_vh_c(D_v),\label{F:KH}
\end{equation}
where $(j)_2=j(j-1)$, $h_c(j)=j(1-a_c)^{j-1}$, and $\E K=n_e(n_e-1)/M$. The quantity $Y_K$ is the
standardization of the occupancy energy, and it is converted back to $X_+$ in \eqref{F:affine}.

After revealing the first $l$ endpoints, denote the occupancy count by $D_v^{(l)}$, and set $m=n_e-l-1$. For
a function $f$ define $\nu^f_m(j)=\E[\Delta f(j+\mathrm{Bin}(m,1/M))]$, where $\Delta f(j)=f(j+1)-f(j)$.
When the next endpoint $V$ lands, the Doob increment of $\sum_vf(D_v)$ is
\begin{equation}
\Delta_f=\nu^f_m(D_V^{(l)})-\frac1M\sum_v\nu^f_m(D_v^{(l)}).\label{F:doob}
\end{equation}
Taking $f(j)=(j)_2$ gives $\nu^f_m(j)=2j+2m/M$, so the Doob increment of the energy is exactly
$2(D_V^{(l)}-l/M)$. The total number of endpoints is fixed sample by sample, and here no Poisson family
replaces the true initial state.

\begin{lemma}[Mean of the bracket]\label{F:lem:bracket}
Denote the matrix bracket of the initial vector $(Y_K,(H-\E H)/(\mu_c\sigma_k\sqrt M))$ by
$(V_{EE},V_{EH},V_{HH})$. Then
\begin{equation}
\E V_{EE}=(1-1/n_e)(1-1/M),\qquad \E V_{EH}=\eta_k+O\bigl((1+|r_M|)/\sqrt M\bigr),\label{F:bracketmean}
\end{equation}
and $\E V_{HH}$ is the same as the initial bracket of Appendix C.
\end{lemma}

\begin{proof}
Expanding by the collision indicators, $\Var K=2n_e(n_e-1)(M-1)/M^2$. For any $h$, with
$B\sim\mathrm{Bin}(n_e-2,1/M)$, there is the finite identity
\begin{equation}
\Cov(K,H)=\frac{n_e(n_e-1)(M-1)}{M^2}\E\Delta^2h(B).\label{F:covKH}
\end{equation}
The proof is as follows. Write $K=\sum_{i\ne j}\mathbf1_{\{V_i=V_j\}}$. Conditioned on $V_1=V_2$, the mean of
$H$ is $\E h(B+2)+(M-1)\E h(B)$; the unconditional mean is
$M\E h(B)+2\E\Delta h(B)+M^{-1}\E\Delta^2h(B)$. The difference is $(1-M^{-1})\E\Delta^2h(B)$, and multiplying
by $n_e(n_e-1)/M$ gives the result.

For $h=h_c$, the Poisson limit gives $\E_{\Pois(\mu)}\Delta^2h_c(D)=a_ce^{-\mu a_c}(\mu a_c-2)$. The function
$\Delta^2h_c$ is bounded, the total-variation error between the binomial and the Poisson is $O(M^{-1})$, the
parameter $\mu$ lies in a compact interval, and the dependence is smooth. By the Doob isometry, $\E V_{EE}$
and $\E V_{EH}$ equal the corresponding variance and covariance divided by the normalizing constants;
substituting $\mu=\mu_c+O(|r_M|/\sqrt M)$ gives \eqref{F:bracketmean}. The limit of the cross term is exactly
$\eta_k$; it comes from the true occupancy model, with no a priori independence between the dynamic noise and
the energy.
\end{proof}

The maximum occupancy count may reach order $\log M$, so the energy increment is unbounded. To obtain
concentration, take $L=\lceil\log M\rceil$ and define the truncated energy $e_L$ by $e_L(0)=0$ and
$\Delta e_L(j)=2\min(j,L)$. For $j\le L+1$ we have $e_L(j)=(j)_2$, and $|\Delta e_L|\le2L$,
$|\Delta^2e_L|\le2$. The true Doob martingale is reconstructed with $e_L$ rather than by simply dropping the
large increments of the original martingale.

\begin{lemma}[Super-polynomial concentration of the truncated bracket]\label{F:lem:conc}
For $(f,g)\in\{(e_L,e_L),(e_L,h_c),(h_c,h_c)\}$,
\begin{equation}
\Pr\Bigl(|V^{(L)}_{fg}-\E V^{(L)}_{fg}|>CL^2(\log M)^2/\sqrt M\Bigr)\le2e^{-c(\log M)^4},\label{F:conc}
\end{equation}
and the difference between $\E V^{(L)}_{fg}$ and the mean of the untruncated bracket, multiplied by any fixed
power of $M$, still tends to zero.
\end{lemma}

\begin{proof}
Changing one input endpoint changes only two occupancy coordinates at each later $l$, each by $1$. For $e_L$,
$\nu^{e_L}_m$ is bounded by $O(L)$, and the per-coordinate change is $O(1)$; for $h_c$, both are $O(1)$. Hence
each $q_l^{fg}=M^{-1}\sum_v\nu^f_m\nu^g_m-\bar\nu^f_m\bar\nu^g_m$ changes by at most $CL^2/M$; after
normalizing and summing over $O(M)$ steps, the sensitivity of the entire bracket to each input coordinate is
at most $CL^2/M$. McDiarmid's inequality gives \eqref{F:conc}.

The probability that the maximum occupancy count exceeds $L$ is at most $M(eC/L)^L$, which is smaller than any
fixed negative power of $M$. On its complement, $K$ equals the truncated energy, while the absolute values of
the two over the whole sample are at most $O(M^2)$. Hence the differences of the center, the variance, and the
covariance with $H$, multiplied by any fixed power of $M$, still tend to zero. Here only the replacement of the
mean is needed; the truncated bracket is not required to equal the original bracket pathwise.
\end{proof}

\subsection{Concatenation and the deterministic trace upper bound}

After the initial state, the labeled stubs are assigned independently and uniformly to the edges, and the
dynamic filtration follows the construction of Appendix C, without revealing the complete grouping in advance.
The initial increments of the two-dimensional vector are
$(\Delta e_L/(\sqrt2\mu\sqrt M),\Delta h_c/(\mu_c\sigma_k\sqrt M))$, and the dynamic increments are
$(0,\text{score martingale increment})$. In the dynamic phase the first coordinate is zero, so it contributes
no new $EE$ or $EH$ bracket; this is an identity on the filtration and does not require an independence
assumption.

By Lemma~\ref{F:lem:bracket}, Lemma~\ref{F:lem:conc}, and input (C2), except for a super-polynomially small
exceptional set, the total matrix bracket $\mathcal V$ satisfies
\begin{equation}
\|\mathcal V-\Sigma\|\le v_M,\qquad \Sigma=\begin{pmatrix}1&\eta_k\\\eta_k&1\end{pmatrix}.\label{F:Sigma}
\end{equation}
$\Sigma$ is positive definite: the initial energy variance is $1$, and the score has strictly positive dynamic
variance $\tau_k^2$, so $|\eta_k|<1$.

After truncation the increment bound is $\beta_M=CL/\sqrt M$, but over the whole sample the trace bracket is
only known to be $O(L^2)$. To perform the exponential tilt, stop the entire vector at the first time the trace
bracket exceeds a fixed constant $C_0>\operatorname{tr}\Sigma+1$. The bracket increment is a positive
semidefinite matrix, so the trace is monotone; the stopping changes the terminal value only when the total
trace exceeds $C_0$, an event of super-polynomially small probability. After stopping,
$\operatorname{tr}\mathcal V\le C_0+C\beta_M^2$ holds deterministically. Below, $(Y^*,S^*)$ denotes this
stopped vector. On the good set it coincides with the original score and the original energy, up to a
super-polynomially small deterministic center term. The subsequent exponential tilt acts only on $S^*$.

\subsection{The two-dimensional rectangle error under a single-coordinate tilt}

This section gives the concrete form of the ``explicit polynomial rate'' of Appendix E of the main text.

\begin{lemma}[Single-coordinate tilt lemma]\label{F:lem:tilt}
Let $\zeta_j=(u_j,v_j)$ be two-dimensional martingale differences with $|\zeta_j|\le\beta$, suppose the trace
of the total matrix bracket $\mathcal V$ has a deterministic upper bound, and suppose
$\Pr(\|\mathcal V-\Sigma\|>v)\le q$, where $\beta,v$ are polynomially small in $M$ and $q$ is
super-polynomially small. Write $(U,S)=\sum_j\zeta_j$, let $1\le t\le C\sqrt{\log M}$, and set the tilted
measure $dP_t/dP=e^{-tS}/\E e^{-tS}$. Then the distance $\epsilon_M$ between the lower-rectangle CDF of
$(U+\eta_kt,S+t)$ under $P_t$ and that of $N(0,\Sigma)$ satisfies
\begin{equation}
\epsilon_M\le C\Bigl\{(\beta+v+t\beta+q_Q)^{1/64}+tv+t^2\beta+t^2v+t^3\beta+q_Q\Bigr\},\label{F:rate}
\end{equation}
where $q_Q$ is super-polynomially small. Substituting $\beta=O(\log M/\sqrt M)$ and \eqref{F:Sigma}, the
right-hand side is a fixed negative power times a polylogarithm, so $t\epsilon_M\to0$.
\end{lemma}

\begin{proof}
\emph{Step 1: product tilt.} Define the measure $Q$ with density
$\prod_je^{-tv_j}/\E(e^{-tv_j}\mid\mathcal F_{j-1})$. The deterministic trace upper bound gives
$\E(dQ/dP)^2\le e^{Ct^2}$, so the bad set is still super-polynomially small under $Q$; the total-variation
distance between the true tilt $P_t$ and $Q$ is $O(t^2v+t^3\beta)$ plus a super-polynomially small quantity.

\emph{Step 2: drift and bracket.} Write $\mathcal B_j=\E(\zeta_j\zeta_j^\top\mid\mathcal F_{j-1})$ and
$e_2=(0,1)^\top$. A stepwise Taylor expansion gives
\begin{equation}
\E_Q(\zeta_j\mid\mathcal F_{j-1})=-t\mathcal B_je_2+\varrho_j,\qquad|\varrho_j|\le Ct^2\beta\operatorname{tr}\mathcal B_j,\label{F:drift}
\end{equation}
and the centered conditional covariance differs from $\mathcal B_j$ in operator norm by at most
$Ct\beta\operatorname{tr}\mathcal B_j$. Let $T$ be the centered vector martingale under $Q$; then on the good
set
\[
\|\langle T\rangle-\Sigma\|\le v+Ct\beta,\qquad|(U+\eta_kt,S+t)-T|\le C(tv+t^2\beta).
\]
The drift $-\eta_kt$ of the first coordinate comes from the cross entries of $\mathcal B_j$ in
\eqref{F:drift}; knowing only that the scalar variance is close to $1$ does not yield this drift, so the $EH$
bracket must be controlled simultaneously.

\emph{Step 3: quantitative two-dimensional rectangle CLT.} For a centered two-dimensional martingale $T$,
suppose the increment bound is $b$, the trace has a deterministic upper bound, and the matrix bracket is within
$w$ of $\Sigma$ with probability at least $1-q'$. Set $\delta=b+w+q'$. For each unit direction $\theta$,
standardize $\theta\cdot T$ by $(\theta^\top\Sigma\theta)^{1/2}$ and stop when the directional bracket first
crosses $1+Cw$; the stopping changes only samples of probability at most $q'$. A one-step Taylor expansion of
the compensated characteristic function gives, direction by direction, for $r=|u|\le\delta^{-1/16}$,
\[
|\varphi_T(u)-e^{-u^\top\Sigma u/2}|\le C\{br^3+b^2r^4+(w+b^2+q')r^2+q'\}\le C\delta(1+r^4),
\]
with the constant uniform in the direction. Here only the characteristic function of the original $T$ is
estimated direction by direction; no common stopping vector is constructed.

Add an independent $N(0,h^2I_2)$ to both sides, with $h=\delta^{1/32}$. Fourier inversion gives a supremum
bound on the difference of the smoothed densities: the integral over the low-frequency ball
$|u|\le\delta^{-1/16}$ is at most $C\delta^{5/8}$, and the high-frequency part is controlled by
$e^{-h^2|u|^2/2}$ and is super-polynomially smaller than $\delta$. The second moments of both sides are
uniformly bounded; restricting the lower rectangle to $[-R,R]^2$ with $R=\delta^{-1/32}$, the probability
outside the region is at most $C/R^2$, and the integral of the density error inside the region is at most
$CR^2\delta^{5/8}$, so the smoothed rectangle error is at most $C\delta^{1/16}$. Finally remove the smoothing:
take the coordinate tolerance $a=h\sqrt{C\log(1/\delta)}$; the probability that the noise exceeds the bound is
smaller than $\delta^2$; the marginal density of the smoothing Gaussian is uniformly bounded, and the
probability of the strip along the rectangle boundary is $O(a)$. Sandwiching the original rectangle between
the $\pm a$ translates of the smoothed rectangle gives
\begin{equation}
d_{\rm rect}(T,N(0,\Sigma))\le C\delta^{1/64}.\label{F:rect}
\end{equation}
This conclusion also covers unbounded endpoints, namely the marginal distributions.

Applying \eqref{F:rect} to the $T$ of Step 2 (with $b=2\beta$ and $w=v+Ct\beta$), controlling the shift
$C(tv+t^2\beta)$ by the anti-concentration of the Gaussian marginals, and adding the total-variation error of
Step 1 gives \eqref{F:rate}.
\end{proof}

\begin{remark}
The rate in \eqref{F:rate} is very weak; it serves only to guarantee $t\epsilon_M\to0$, and is not used to
predict the error at finite $M$.
\end{remark}

\subsection{The macroscopic conditional CDF of the all-positive-sign measurement}

\begin{proposition}[Macroscopic conditional law for all-positive signs]\label{F:prop:cdf}
Under \eqref{F:range}, with the cutoff of Lemma~\ref{F:lem:cutoff}, the Plain model satisfies
\begin{equation}
\sup_x\left|\Pr(X_+\le x\mid A)-\frac{\Phi_2(x,-t_0;\eta_k)}{\Phi(-t_0)}\right|\to0.\label{F:cdf}
\end{equation}
Equivalently, $X_++\eta_kt_0$ conditioned on $A$ tends to $N(0,1-\eta_k^2)$.
\end{proposition}

\begin{proof}
\emph{Inverse tilt.} The stopped $S^*$ satisfies the scalar Laplace approximation
$\log\E e^{-tS^*}=t^2/2+o(1)$; take $t=b_M$. For any $y$, set
$\mathcal K_M(y)=\E_t[e^{t(S^*+t)}\mathbf1_{\{S^*+t\le0,\ Y^*+\eta_kt\le y\}}]$. Integration by parts in the
Lebesgue--Stieltjes sense writes $\mathcal K_M$ as an integral of the tilted joint CDF, so
Lemma~\ref{F:lem:tilt} gives $\sup_y|\mathcal K_M(y)-\mathcal K_G(y)|\le2\epsilon_M$, where $\mathcal K_G$
corresponds to the target Gaussian. We have
$\mathcal K_G(\infty)=e^{t^2/2}\Phi(-t)\sim(\sqrt{2\pi}t)^{-1}$. The exact inverse-tilt formula is
\[
\Pr(Y^*\le x,\ S^*\le-t)=\E e^{-tS^*}\,e^{-t^2}\,\mathcal K_M(x+\eta_kt),
\]
Controlling the numerator and denominator simultaneously by the difference between $\mathcal K_M$ and
$\mathcal K_G$, the error of the conditional CDF is at most
$O(\epsilon_M/\mathcal K_G(\infty))=O(t\epsilon_M)=o(1)$:
\[
\sup_x\left|\Pr(Y^*\le x\mid S^*\le-t)-\frac{\Phi_2(x,-t;\eta_k)}{\Phi(-t)}\right|\to0.
\]

\emph{Transfer to the macroscopic event.} By input (C3) and the truncation and stopping bad sets of this
appendix, $\Pr(A\triangle\{S^*\le-t\})=o(\Phi(-t))$. This probability bound controls simultaneously the joint
sub-probability error of all $\{Y_K\le x\}$, without an energy-weighted bound, so the conditional CDF can be
transferred. On the good set $Y_K$ and $Y^*$ differ only by a super-polynomially small amount; the density of
the selection Gaussian is uniformly bounded, and, given the score, the conditional density of the energy does
not exceed $[2\pi(1-\eta_k^2)]^{-1/2}$, and this remains true after mixing. Replace $t$ by $t_0$: the two
Gaussian-tail conditioning events are nested, and the total-variation distance of the conditional distributions
is at most the relative difference of the tail probabilities $O(t_0|t-t_0|)=o(1)$. This yields the conclusion
for $Y_K$ in the configuration model.

\emph{Fixed planting.} Fix $l$ special edges, construct the initial two-dimensional martingale with $k(d-l)$
regular endpoints, and then append the special-edge filtration of Appendix C. Replacing $\mu$ by $k(d-l)/M$
changes only the $O_l(M^{-1})$ error in \eqref{F:bracketmean}. When the maximum regular degree does not exceed
$L$, the $kl$ special incidences change $K$ by at most $C_l(L+1)$ and $H$ by $C_l$; after normalization this
is $O_l(L/\sqrt M)$, and putting it into the perturbation of the energy coordinate and the bandwidth of the
macroscopic event, multiplied by $t$, still tends to zero. Write
$G_M(x)=\Phi_2(x,-t_0;\eta_k)/\Phi(-t_0)$; then for each fixed $l$,
\begin{equation}
\sup_x\left|\Pr\nolimits_l(Y_K\le x,\ A)/p_M-G_M(x)\right|\to0.\label{F:plant}
\end{equation}
This formula already contains the relative precision of the probability denominator, which an ordinary
finite-planting CLT cannot replace.

\emph{Plain transfer.} Let $N^{\rm bad}_M$ be the number of bad edges, with $dp_M^{\rm bad}\to\Lambda_k$.
Specifying $l$ bad edges is exactly the allowed special-edge model, so for each fixed $l$,
\[
\sup_x\left|\E[\mathbf1_{\{Y_K\le x,\ A\}}(N^{\rm bad}_M)_l]/p_M-\Lambda_k^lG_M(x)\right|\to0.
\]
Apply finite-order Bonferroni upper and lower bounds to the nonnegative indicators. Fix the truncation order
and let $M\to\infty$; the finite-term error is uniform in $x$. The difference between the two bounding sums and
$e^{-\Lambda_k}$, multiplied by $G_M(x)\le1$, then vanishes as the truncation order increases. This gives
$\sup_x|\Pr(Y_K\le x,A,N^{\rm bad}_M=0)/p_M-e^{-\Lambda_k}G_M(x)|\to0$, and dividing by
$\Pr(A,N^{\rm bad}_M=0)/p_M\to e^{-\Lambda_k}$ yields the conditional CDF of the Plain model.

\emph{Conversion back to $X_+$.} In the Plain model exactly
\begin{equation}
X_+=\frac{\sqrt2\mu\sqrt M}{k(1-k/M)s_M}Y_K+\frac{(k-1)\mu}{k(1-k/M)s_M},\label{F:affine}
\end{equation}
that is, $X_+=(1+O(M^{-1}))Y_K+O(M^{-1/2})$. The quantity $Y_K+\eta_kt_0$ is tight under $A$, so the
transformation error is $o_P(1)$, and using the density upper bound above gives \eqref{F:cdf}. The energy of
the all-negative-sign case equals that of the all-positive-sign case.
\end{proof}

\subsection{Conditional moments and uniform integrability}

Convergence in distribution does not yield convergence of moments after shifting the center. This section
controls the tails directly under the macroscopic conditioning.

\begin{proposition}[Macroscopic conditional moments]\label{F:prop:moments}
Under the hypotheses of Proposition~\ref{F:prop:cdf}, the moments of every fixed positive integer order and the
absolute moments of $W^+_M=X_++\eta_kt_0$ under $A$ converge to the corresponding moments of
$N(0,1-\eta_k^2)$. In particular,
\begin{equation}
\E[X_+\mid A]=-\eta_k\ell(t_0)+o(1),\qquad\Var(X_+\mid A)=1-\eta_k^2+o(1).\label{F:moments}
\end{equation}
For any fixed sign composition independent of the mapping,
$\E[X_s\mid A]=\gamma_M\bigl[-\eta_k\ell(t_0)+o(1)\bigr]$, where the $o(1)$ is the same all-positive-sign
remainder, hence uniform in the composition, and no convergence of $\rho_M$ is required.
\end{proposition}

\begin{proof}
\emph{Quantitative CDF input.} We first work in the configuration model. The polynomial rectangle error of
Lemma~\ref{F:lem:tilt}, the factor $t$ of the inverse tilt, the relative bandwidth error
$O(t(\delta_M+\epsilon_M))$, and the super-polynomially small bad set divided by $p_M$, together give
\begin{equation}
d_K\bigl(\mathcal L(Y_K\mid A),\ \mathcal G_t\bigr)\le e_M=O(M^{-c}(\log M)^b),\label{F:dK}
\end{equation}
where $\mathcal G_t$ is the distribution of the first coordinate of a standard two-dimensional Gaussian with
correlation $\eta_k$, conditioned on the second coordinate being $\le-t$, with $c>0$ and $b<\infty$ fixed. The
Bonferroni bounds have not been used here; when transferring to the Plain model later, only qualitative CDF
convergence is needed.

\emph{Tilted tail bound over the whole sample.} Continue with the product tilt $Q$ of \eqref{F:drift}. Under
$Q$, $T_1=\sum_j(u_j-\E_Q[u_j\mid\mathcal F_{j-1}])$ is a centered martingale. By \eqref{F:drift}, over the
whole sample,
\[
\Bigl|\sum_j\E_Q[u_j\mid\mathcal F_{j-1}]\Bigr|\le t\sum_j|(\mathcal B_j)_{12}|+Ct^2\beta\sum_j\operatorname{tr}\mathcal B_j\le C_1t,
\]
Here only $|(\mathcal B_j)_{12}|\le\operatorname{tr}\mathcal B_j/2$, the deterministic trace upper bound, and
$t\beta\to0$ are used; the cross bracket on the bad set is not required to be close to $\eta_k$. The increments
of $T_1$ are at most $2\beta$ and the sum of the predictable variances is at most $C_2$, so the
Bernstein--Freedman inequality gives $Q(|T_1|>x)\le2\exp\{-x^2/[2(C_2+2\beta x/3)]\}$. Set $W^*=Y^*+\eta_kt$;
then $|W^*-T_1|\le C_3t$, so for $x\ge2C_3t$,
\begin{equation}
Q(|W^*|>x)\le2\exp\{-cx^2/(1+\beta x)\}.\label{F:Qtail}
\end{equation}

\emph{Unbounded moments without total variation.} The density of the true tilt relative to $Q$ is exactly
$\prod_jm_j(-t)/\E e^{-tS^*}$, where $m_j(-t)=\E(e^{-tv_j}\mid\mathcal F_{j-1})$. The deterministic trace
upper bound and $t\beta\to0$ give $\prod_jm_j\le e^{Ct^2}$; combined with $\E e^{-tS^*}\sim e^{t^2/2}$, this
yields the deterministic bound $dP_t/dQ\le e^{C_4t^2}\le M^{D_0}$ over the whole sample, with $D_0$ fixed.
This crude bound is used together with martingale tail bounds whose order can be raised arbitrarily;
multiplying a small total-variation distance by a growing energy maximum does not yield a vanishing error.

\emph{Truncated tails.} Fix a positive integer $q$. The actual energy satisfies $|Y_K|\le CM^{3/2}$, and the
stopped martingale has at most $O(M)$ steps, each of size $O(\log M/\sqrt M)$, so $|W^*|\le CM^2$ is a
deterministic crude bound. Take $R_M=A_q\sqrt{\log M}$ with $A_q$ large enough that $R_M\ge2C_3t$ and that
\eqref{F:Qtail} gives $Q(|W^*|>R_M)\le2M^{-H_q}$ with $H_q>D_0+2q+2$. Outside $\mathcal B_M$, the event $A$
implies $S^*\le-t+C\delta_M$; for any nonnegative $f(W^*)$, the inverse tilt gives
\[
\E[f(W^*);A\setminus\mathcal B_M]\le\E e^{-tS^*}e^{-t^2}e^{Ct\delta_M}\E_tf(W^*)\le Cp_M\,t\,M^{D_0}\E_Qf(W^*),
\]
where the last step uses $\E e^{-tS^*}e^{-t^2}/p_M=O(t)$. Taking
$f=|W^*|^q\mathbf1_{\{|W^*|>R_M\}}$ gives
$\E[|W^*|^q\mathbf1_{\{|W^*|>R_M\}};A\setminus\mathcal B_M]/p_M\le CtM^{D_0+2q-H_q}=o(1)$. The contribution of
the bad set is at most $CM^{2q}\Pr(\mathcal B_M)/p_M=o(1)$. Hence
\begin{equation}
\E[|W_K|^q\mathbf1_{\{|W_K|>R_M\}}\mid A]\to0,\qquad W_K=Y_K+\eta_kt.\label{F:tailcut}
\end{equation}
Since $R_M\to\infty$, \eqref{F:tailcut} by itself is not yet uniform integrability.

\emph{Between a fixed threshold and $R_M$.} Let $Z_0,G$ be independent standard normals; the shifted-center
variable of the selection Gaussian can be written as
$W^G_t=\eta_k(Z_0+t)+\sqrt{1-\eta_k^2}\,G$, conditioned on $Z_0\le-t$. For $t\ge1$ and $x\ge0$,
\[
\Pr(-(Z_0+t)>x\mid Z_0\le-t)=\frac{\Phi(-(t+x))}{\Phi(-t)}\le e^{-tx-x^2/2},
\]
which follows by shifting the integration variable of the numerator by $x$. Hence the moments of every order of
the overshoot tend to zero, and the moments of every order of $W^G_t$ tend to the corresponding moments of
$N(0,1-\eta_k^2)$. For the $q$-th moment use $\mathrm{clip}(w,-R_M,R_M)^q$, and for the absolute moment use
$\min(|w|,R_M)^q$; their total variation is at most $2R_M^q$, and Stieltjes integration by parts with
\eqref{F:dK} gives that the difference of the truncated expectations is at most $2R_M^qe_M=o(1)$. Then
\eqref{F:tailcut} and the uniform tail control of the Gaussian remove the truncation, giving the convergence of
all moments in the configuration model. Taking the conclusion of order $q+1$ gives the uniform integrability
\begin{equation}
\sup_{M\text{ sufficiently large}}\E[|W_K|^q\mathbf1_{\{|W_K|>R\}}\mid A]\le C_q/R.\label{F:UI}
\end{equation}

\emph{Plain transfer.} The scalar relative Bonferroni already gives
$\Pr_{\rm conf}(A,N^{\rm bad}_M=0)/\Pr_{\rm conf}(A)\to e^{-\Lambda_k}>0$. For any nonnegative $f$,
\[
\E_{\rm Plain}[f(Y_K)\mid A]=\frac{\E_{\rm conf}[f(Y_K)\mathbf1_{\{N^{\rm bad}_M=0\}}\mid A]}{\Pr_{\rm conf}(N^{\rm bad}_M=0\mid A)}\le C\,\E_{\rm conf}[f(Y_K)\mid A],
\]
so \eqref{F:UI} carries over directly to the Plain model. This inequality transfers only the tail control; the
limiting distribution comes from the Plain CDF transfer of Proposition~\ref{F:prop:cdf}. Combining the two
gives all moment limits in the Plain model. Finally, convert back to $X_+$ by \eqref{F:affine}:
$X_++\eta_kt_0=(1+O(M^{-1}))(Y_K+\eta_kt)+O(M^{-1/2})+\eta_k(t_0-t)+O(t/M)$, and the last several terms tend
to zero.

\emph{Mean and variance.} The first- and second-moment conclusions give $\E[X_+\mid A]=-\eta_kt_0+o(1)$ and
$\Var(X_+\mid A)=1-\eta_k^2+o(1)$. The Mills ratio satisfies $\ell(t)-t=O(1/t)$, so the mean can also be
written in the form of \eqref{F:moments}. This is only a rewriting; the remainder is still $o(1)$, and it does
not identify the coefficient of order $1/t_0$.

\emph{Mean of an arbitrary composition.} The event $A$ depends only on the unsigned mapping and is invariant
under permutation of the element labels, so the exact projection of Eq.~(7) of the main text gives
$\E[X_s\mid A]=\gamma_M\E[X_+\mid A]$. This step does not require a sign CLT under a rare event.
\end{proof}

\subsection{Arbitrary sign composition: conditional CDF, variance, and endpoints}

\begin{proposition}[Macroscopic conditional law for an arbitrary composition]\label{F:prop:sign}
Under the hypotheses of Proposition~\ref{F:prop:cdf}, with an arbitrary fixed sign composition, $\rho_M$ need
not converge. Set $G_{M,s}(x)=\Phi_2(x,-t_0;\eta_k\rho_M^2)/\Phi(-t_0)$; then
\begin{equation}
\sup_x|\Pr(X_s\le x\mid A)-G_{M,s}(x)|\to0,\qquad\Var(X_s\mid A)=1-\eta_k^2\rho_M^4+o(1).\label{F:sign}
\end{equation}
\end{proposition}

The idea of the proof is to avoid redoing a signed tilt for each sign composition. First generate the unsigned
Plain hypergraph, then assign the given numbers of positive and negative signs uniformly; this has the same
joint law as the original model for $(X_s,A)$, and $A$ selects only the graph. Hence only two things are
needed: the required graph properties still hold with probability tending to one under $A$; and on these fixed
good graphs, the sign CLT holds uniformly. The graph properties used in Theorem 2 of the main text hold only to
precision $o_P(1)$, so they cannot be divided directly by the rare $\Pr(A)$. Below, the exceptional probability
is strengthened to be super-polynomially small.

We use the overlap-graph notation of Section 5.6: $m_{ij}=|S_i\cap S_j|$, $\Omega=\sum_{i<j}m_{ij}$,
$\Omega_2=\sum_{i<j}m_{ij}^2$, $\ell_i=\sum_{j\ne i}m_{ij}$, $\Omega_\ell=\sum_i\ell_i^2$, and
$\bar\ell=2\Omega/d$.

\begin{lemma}[Super-polynomial concentration of the overlap-graph statistics]\label{F:lem:graph}
There is a deterministic set of graphs $\mathcal H_M$ such that $\Pr(\mathcal G_M\notin\mathcal H_M)$ is
super-polynomially small, and all graphs in $\mathcal H_M$ satisfy uniformly
\begin{equation}
\frac\Omega M=\frac{\mu_c^2}2+o(1),\quad\frac{\Omega_2}M=\frac{\mu_c^2}2+o(1),\quad\frac1M\sum_i(\ell_i-\bar\ell)^2=\mu_c^2+o(1),\quad\max_i\ell_i\le kL.\label{F:goodgraph}
\end{equation}
All the $o(1)$ terms can be taken as $O((\log M)^4/\sqrt M+|r_M|/\sqrt M)$. In particular,
$\Pr(\mathcal G_M\notin\mathcal H_M\mid A)\to0$, and this exceptional probability, multiplied by any fixed
power of $M$ and divided by $\Pr(A)$, still tends to zero.
\end{lemma}

\begin{proof}
Let $L=\lceil\log M\rceil$ and $\mathcal D_L=\{\max_vD_v\le L\}$. Since $D_v\sim\mathrm{Bin}(d,k/M)$, we have
$\Pr(\mathcal D_L^c)\le M(eC/L)^L$, which is super-polynomially small.

\emph{Lipschitz bounds.} For two graphs in $\mathcal D_L$, successively replace the edges with different labels
by the edges of the other graph. The incidence set of every intermediate graph is contained in the union of
the two endpoint graphs, so the vertex degrees do not exceed $2L$. Replacing one edge in such an intermediate
graph affects only the $O(L)$ edge labels overlapping the old or the new edge, and each $m_{ij}\le k$, so the
changes of $\Omega$ and $\Omega_2$ are both $O(L)$. For the replaced edge $\ell_i=O(L)$, so the squared change
is $O(L^2)$; for the remaining $O(L)$ affected labels, $\ell_j$ changes by $O(1)$, so each square changes by
$O(L)$. Hence the total change of $\Omega_\ell$ is $O(L^2)$. Thus the three statistics are $CL^2$-Lipschitz on
$\mathcal D_L$ with respect to the Hamming distance on the edge labels. Here the global Lipschitz property is
proved along a path on the union of two good graphs, without assuming that the one-step path stays inside
$\mathcal D_L$.

\emph{Extension and concentration.} For each statistic $f$ take the McShane extension
$\widetilde f(\mathcal G)=\inf_{\mathcal G'\in\mathcal D_L}\{f(\mathcal G')+CL^2d_H(\mathcal G,\mathcal G')\}$
and clip it to the deterministic nonnegative polynomial range of $f$. The extension equals $f$ on the good set
and remains $CL^2$-Lipschitz; the bad-set probability is super-polynomially small, and both $f$ and the clipped
extension are at most $O(M^3)$, so $|\E f-\E\widetilde f|$ is super-polynomially small. Applying McDiarmid's
inequality to the independent edges,
\[
\Pr\bigl(|f-\E f|>C\sqrt M(\log M)^4\bigr)\le\Pr(\mathcal D_L^c)+2e^{-c(\log M)^4},\qquad f=\Omega,\Omega_2,\Omega_\ell.
\]

\emph{Deterministic means.} For the overlap $H$ of two independent uniform $k$-subsets, write
$h_1=\E H=k^2/M$ and $h_2=\E H^2=k^2/M+k^2(k-1)^2/[M(M-1)]$. Exactly,
\[
\E\Omega=\tbinom d2h_1,\qquad\E\Omega_2=\tbinom d2h_2,\qquad\E\Omega_\ell=d\{(d-1)h_2+(d-1)(d-2)h_1^2\}.
\]
The last formula holds because, given $S_i$, the $m_{ij}$ are independent and their conditional first and
second moments do not depend on $S_i$. Taking the good set as the intersection of $\mathcal D_L$ with the three
deviation-threshold events gives \eqref{F:goodgraph}; the third relation is computed via
$\Omega_\ell-4\Omega^2/d$, in which the deterministic constants of the two terms cancel. By input (C4),
$\Pr(A)\sim p_M$ is no smaller than a fixed negative power divided by a polylogarithm, so the conditional
exceptional probability tends to zero.
\end{proof}

\begin{lemma}[Uniform sign CLT on the good graphs]\label{F:lem:unifsign}
Take the exact decomposition $X_s=\gamma_MX_++V_M$ with $\E_s[V_M\mid\mathcal G_M]=0$. If $\rho_M\to\rho$ and
$|\rho|<1$, then for each fixed $u$,
\begin{equation}
\sup_{\mathcal G\in\mathcal H_M}\Bigl|\E_s[e^{iuV_M}\mid\mathcal G]-e^{-(1-\rho^4)u^2/2}\Bigr|\to0.\label{F:unifsign}
\end{equation}
\end{lemma}

\begin{proof}
The sign argument of Section 5.6 and Appendix A is an argument over a fixed sequence of graphs; it uses only
\eqref{F:goodgraph}, $\Omega=O(M)$, $\Omega_\ell=O(M)$, and the condition that the positive-to-negative ratio
stays at a fixed distance from the endpoints. Specifically, the within-group variance bound of the batch flip
uses only $\Omega_\ell=O(M)$, and the variance of the flipped subset and the expectations of the internal
edges use only $\Omega_\ell$ and $\Omega$; the number of local terms in the main term is $O(M)$, the maximum
degree of the dependency graph and the magnitude of the terms are both $O(\log M)$, and the normalized
connected cumulants of fixed order $r\ge3$ are $O_r(M^{1-r/2}\,\mathrm{polylog}\,M)$; the variance of the main
term is determined by \eqref{F:goodgraph}. Hence for any sequence $\mathcal G_M\in\mathcal H_M$,
$\mathcal L_s(V_M\mid\mathcal G_M)\Rightarrow N(0,1-\rho^4)$. If \eqref{F:unifsign} failed, one could extract
along a subsequence a deterministic graph violating it; that sequence of graphs would still satisfy all the
deterministic bounds, contradicting the conclusion above.
\end{proof}

\begin{lemma}[Per-graph conditional variance and the endpoints]\label{F:lem:endpoint}
Uniformly over all feasible compositions, on $\mathcal H_M$,
\begin{equation}
\Var_s(V_M\mid\mathcal G)=1-\rho_M^4+o(1).\label{F:pergraphvar}
\end{equation}
\end{lemma}

\begin{proof}
One cannot first compute the unconditional $\E V_M^2$ and then divide by the rare $\Pr(A)$, so use instead the
exact per-graph variance. Let $a_4$ be the mean of the product of the signs of four distinct labels. The
finite identity gives
\[
\Var_s\Bigl(\sum_{i<j}m_{ij}s_is_j\Bigm|\mathcal G\Bigr)=\Omega_2(1-2\gamma_M+a_4)+\Omega_\ell(\gamma_M-a_4)+\Omega^2(a_4-\gamma_M^2).
\]
Set $v=1-\rho_M^2$ and expand this finite rational expression uniformly over all feasible compositions; this
gives
\[
v^2\Omega_2+\rho_M^2v(\Omega_\ell-4\Omega^2/d)+O(\Omega_2/d+\Omega_\ell/d+\Omega^2/d^2).
\]
On $\mathcal H_M$ we have $\Omega,\Omega_2,\Omega_\ell=O(M)$, so the remainder is $O(1)$; multiplying by the
normalization factor $4/[k^2(1-k/M)^2s_M^2]\sim2/(\mu_c^2M)$ gives \eqref{F:pergraphvar}. This also yields the
crude upper bound $C(v+M^{-1})$. When $|\rho_M|\to1$, \eqref{F:pergraphvar} and the per-graph Chebyshev
inequality give $V_M\to0$, uniformly over the good graphs; this covers any rate of approach to the endpoints.
\end{proof}

\begin{proof}[Proof of Proposition~\ref{F:prop:sign}]
Let $U_M=X_++\eta_kt_0$. Proposition~\ref{F:prop:cdf} gives $U_M\mid A\Rightarrow N(0,1-\eta_k^2)$. For any
subsequence with $\rho_M\to\rho$ and $|\rho|<1$, first condition on the graph: on the good graphs the inner
error tends to zero uniformly by \eqref{F:unifsign}, on the bad graphs the characteristic function is bounded
by one in absolute value, and by Lemma~\ref{F:lem:graph} the probability of the bad graphs under $A$ tends to
zero. Hence
\[
\E[e^{iaU_M+ibV_M}\mid A]-\E[e^{iaU_M}\mid A]\,e^{-(1-\rho^4)b^2/2}\to0,
\]
that is, $U_M$ and $V_M$ are asymptotically independent under $A$. When $|\rho|=1$, use instead the convergence
in conditional probability of Lemma~\ref{F:lem:endpoint}.

Keeping the actual shifted center:
\[
X_s+\eta_k\rho_M^2t_0=\gamma_MU_M+V_M+\eta_k(\rho_M^2-\gamma_M)t_0,
\]
The last term is $O(t_0/d)\to0$. Here $\rho_M^2t_0$ is not replaced by $\rho^2t_0$ in the absence of a rate.
Hence $\mathcal L(X_s+\eta_k\rho_M^2t_0\mid A)\Rightarrow N(0,1-\eta_k^2\rho^4)$. The selection baseline
$G_{M,s}$ has the same limit under the same shifted center: it can be written as
$c(Z+t_0)+\sqrt{1-c^2}\,G$ conditioned on $Z\le-t_0$, where $c=\eta_k\rho_M^2$, and the Gaussian overshoot
tends to zero. The two distributions tend to the same continuous distribution, so their CDF distance tends to
zero. When $\rho_M$ does not converge, take a subsequence on which the CDF distance has a positive lower bound,
and by the compactness of $[-1,1]$ extract a further subsequence on which $\rho_M$ converges, contradicting the
conclusion above. This proves the first formula of \eqref{F:sign}.

The variance follows from the law of total variance:
\[
\Var(X_s\mid A)=\gamma_M^2\Var(X_+\mid A)+\E[\Var_s(V_M\mid\mathcal G_M)\mid A].
\]
On the good graphs use \eqref{F:pergraphvar}; on the bad graphs use the deterministic polynomial upper bound on
the energy and the super-polynomial precision of Lemma~\ref{F:lem:graph}. Substituting
$\Var(X_+\mid A)\to1-\eta_k^2$ from Proposition~\ref{F:prop:moments} and noting
$\gamma_M^2-\rho_M^4=O(1/d)$ gives the second formula of \eqref{F:sign}. The conditional mean still uses the
exact projection and cannot be replaced by the CDF convergence.
\end{proof}

\begin{remark}[Balancing rate]
In the fixed critical window, $\rho_M\to0$ suffices to recover the standard normal. In the moderate deviation
regime, the center $-\eta_k\rho_M^2t_0$ of the macroscopic layer grows with $t_0$: if $\rho_M^2t_0\to0$ (which
the exactly balanced composition satisfies), then $X_s\mid A\Rightarrow N(0,1)$; if $\rho_M^2t_0\to
b\in(0,\infty)$, the limit is $N(-\eta_kb,1)$. This refers to the macroscopic layer. When the weight of the
repeated-edge layer in $F$ tends to one, the failure conditional law still recovers normality.
\end{remark}

\subsection{Mixing of the two regimes and the weighted mean}

Lemma~\ref{F:lem:cutoff} gives the exact sizes of the parts in the set decomposition $F=A\cup D\cup R$. The
planting transfer of Appendix B gives $\sup_x|\Pr(X_s\le x\mid D)-\Phi(x)|\to0$, which requires only that
$d/M$ tends to a finite positive number, and holds equally for $\alpha_M\to\alpha_c$. Hence
\[
\Pr(X_s\le x,F)=\Pr(A)\Pr(X_s\le x\mid A)+\lambda_M\Phi(x)+o(\lambda_M),
\]
with the error uniform in $x$. Substituting input (C4) and Proposition~\ref{F:prop:sign} gives the mixture
formula of Proposition 8 of the main text. The argument of Section E.2 of Appendix E of the main text is
exactly this step; here only the cutoff condition is added.

The $o(1)$ error of the mixture CDF cannot be multiplied directly by the growing macroscopic center, so the
conditional mean must be controlled separately on the unnormalized first moment.

\begin{proposition}[Weighted mean]\label{F:prop:wm}
Under \eqref{F:range}, take the all-same-sign measurement $Y=s_MX_+$, let $N$ be the number of repeated edge
pairs, and set $D=\{N\ge1\}$. Then
\begin{equation}
\E[Y;F]=2\lambda_M-s_M\eta_k\phi(t_0)+o\bigl(\lambda_M+s_M\phi(t_0)\bigr).\label{F:wm}
\end{equation}
\end{proposition}

\begin{proof}
The exact density of the random-label planting gives $\E[YN]=2\lambda_M$. Set $R=F\setminus(A\cup D)$; per
graph there is the identity
\begin{equation}
Y\mathbf1_F-Y\mathbf1_A-YN=Y\mathbf1_R-YN\mathbf1_A-Y(N-1)\mathbf1_{D\setminus A}.\label{F:wm-id}
\end{equation}
Below, the absolute first moments of the three terms on the right are estimated one by one.

\emph{Third-kind failure.} The moment estimate of the covering event in Section B.5 of Appendix B uses only
$d\le AM$, so it holds uniformly on the current moving load: $\E[|Y|;R]\le C\sqrt M\,M^{-b_k}=o(\lambda_M)$,
since $b_k-(k-2)\ge1$. Here the core on $R$ has at most $\varepsilon_0M$ edges, and no fixed subcritical gap
is needed.

\emph{Multiplicity of the repeated pairs.} The joint condition of any two specified repeated pairs involves at
most four special labels, and the inner, cross, and outer energy decomposition still gives
$\E|Y|\le C\sqrt M$; also $\E[N(N-1)]\le\lambda_M^2$, so
$\E[|Y|(N-1);D]\le C\sqrt M\lambda_M^2=o(\lambda_M)$.

\emph{Intersection of the macroscopic layer and the repeated layer.} Let $Q$ be the repeated-planting
distribution of the random-label pairs, whose density is exactly $N/\lambda_M$. Then
\[
\E[|Y|N;A]=\lambda_M\E_Q[|Y|;A]\le\lambda_M\sqrt{\E_QY^2}\sqrt{Q(A)}\le C\lambda_M\sqrt{Mp_M}.
\]
The bound $\E_QY^2=O(M)$ comes from the energy decomposition after repeated planting in Appendix B and the
unconditional variance of the outer part. The estimate $Q(A)\sim p_M$ comes from the fixed-planting version of
input (C4): treat two identical legitimate $k$-subsets as finite special edges, then apply the relative
Bonferroni transfer to the remaining regular edges. Set $b=s_Mt_0/\lambda_M$; from
$\sqrt{p}/(1+bp)\le1/(2\sqrt b)$,
\[
\frac{\lambda_M\sqrt{Mp_M}}{\lambda_M+s_Mp_Mt_0}\le\frac12\sqrt{\frac{M\lambda_M}{s_Mt_0}}=O\Bigl(M^{(5/2-k)/2}/\sqrt{t_0}\Bigr)=o(1).
\]
As $t_0\to\infty$, $p_M\ell(t_0)=\phi(t_0)\approx p_Mt_0$, so the intersection term is
$o(\lambda_M+s_M\phi(t_0))$, with no need for a refined weighted local response.

\emph{Macroscopic main term.} By input (C4) and Proposition~\ref{F:prop:moments},
$\E[Y;A]=s_M\Pr(A)[-\eta_k\ell(t_0)+o(1)]=-s_M\eta_k\phi(t_0)+o(s_M\phi(t_0))$. The probability needs only
relative $o(1)$ precision, since the target remainder is $o(\lambda_M+s_M\phi(t_0))$.

Substituting the three estimates and the main term into \eqref{F:wm-id} gives \eqref{F:wm}.
\end{proof}

Combining $\Pr(F)=(p_M+\lambda_M)(1+o(1))$ with $c_F\binom d2=\E[Y\mid F]/2$, \eqref{F:wm} yields the
expansion of $c_F\binom d2$ in Appendix E of the main text, and hence the sign-reversal-layer conclusion of
Proposition 10. For any fixed composition, $\E[\widehat d_M-d;F]$ equals exactly the all-same-sign expectation
multiplied by $\gamma_M$, so both the main term and the remainder project by this factor.

\end{document}
